\documentclass[10pt,a4paper]{amsart}
\usepackage{amsmath,amssymb,amsthm,graphicx}
% Retained minimal preamble for the Stacks Project chapter sites-cohomology.tex (Cohomology on Sites).
% Verbatim excerpt of the chapter preamble: Xy-pic setup (native preamble lines 14-18) and the native
% theorem declarations with the mathematical macros (native preamble lines 154-217). Copied mechanically;
% nothing added, removed or redefined. The wrapper class and the amsmath/amssymb/amsthm/graphicx base
% packages are supplied by the assembler, so the class/externaldocument part of the native preamble is omitted.
% For commutative diagrams we use Xy-pic
\usepackage[all]{xy}

% We use 2cell for 2-commutative diagrams.
\xyoption{2cell}
% Theorem environments.
%
\theoremstyle{plain}
\newtheorem{theorem}[subsection]{Theorem}
\newtheorem{proposition}[subsection]{Proposition}
\newtheorem{lemma}[subsection]{Lemma}

\theoremstyle{definition}
\newtheorem{definition}[subsection]{Definition}
\newtheorem{example}[subsection]{Example}
\newtheorem{exercise}[subsection]{Exercise}
\newtheorem{situation}[subsection]{Situation}

\theoremstyle{remark}
\newtheorem{remark}[subsection]{Remark}
\newtheorem{remarks}[subsection]{Remarks}

\numberwithin{equation}{subsection}

% Macros
%
\def\lim{\mathop{\mathrm{lim}}\nolimits}
\def\colim{\mathop{\mathrm{colim}}\nolimits}
\def\Spec{\mathop{\mathrm{Spec}}}
\def\Hom{\mathop{\mathrm{Hom}}\nolimits}
\def\Ext{\mathop{\mathrm{Ext}}\nolimits}
\def\SheafHom{\mathop{\mathcal{H}\!\mathit{om}}\nolimits}
\def\SheafExt{\mathop{\mathcal{E}\!\mathit{xt}}\nolimits}
\def\Sch{\mathit{Sch}}
\def\Mor{\mathop{\mathrm{Mor}}\nolimits}
\def\Ob{\mathop{\mathrm{Ob}}\nolimits}
\def\Sh{\mathop{\mathit{Sh}}\nolimits}
\def\NL{\mathop{N\!L}\nolimits}
\def\CH{\mathop{\mathrm{CH}}\nolimits}
\def\proetale{{pro\text{-}\acute{e}tale}}
\def\etale{{\acute{e}tale}}
\def\QCoh{\mathit{QCoh}}
\def\Ker{\mathop{\mathrm{Ker}}}
\def\Im{\mathop{\mathrm{Im}}}
\def\Coker{\mathop{\mathrm{Coker}}}
\def\Coim{\mathop{\mathrm{Coim}}}

% Boxtimes
%
\DeclareMathSymbol{\boxtimes}{\mathbin}{AMSa}{"02}

%
% Macros for moduli stacks/spaces
%
\def\QCohstack{\mathcal{QC}\!\mathit{oh}}
\def\Cohstack{\mathcal{C}\!\mathit{oh}}
\def\Spacesstack{\mathcal{S}\!\mathit{paces}}
\def\Quotfunctor{\mathrm{Quot}}
\def\Hilbfunctor{\mathrm{Hilb}}
\def\Curvesstack{\mathcal{C}\!\mathit{urves}}
\def\Polarizedstack{\mathcal{P}\!\mathit{olarized}}
\def\Complexesstack{\mathcal{C}\!\mathit{omplexes}}
% \Pic is the operator that assigns to X its picard group, usage \Pic(X)
% \Picardstack_{X/B} denotes the Picard stack of X over B
% \Picardfunctor_{X/B} denotes the Picard functor of X over B
\def\Pic{\mathop{\mathrm{Pic}}\nolimits}
\def\Picardstack{\mathcal{P}\!\mathit{ic}}
\def\Picardfunctor{\mathrm{Pic}}
\def\Deformationcategory{\mathcal{D}\!\mathit{ef}}

% --- appended 2026-10-12 (Chief): the project own definitions of the reference/slogan/history
% display environments, copied mechanically from the shared Stacks preamble lines 8-12
% (\usepackage{verbatim} plus the three \newenvironment lines, which the chapter preamble keeps
% verbatim from the project preamble). Nothing added, removed or redefined: these environments
% carry no mathematical content and the project itself typesets them as comments. Without these
% three lines a counted statement whose own text uses them fails with Environment slogan undefined.
\usepackage{verbatim}
\newenvironment{reference}{\comment}{\endcomment}
\newenvironment{slogan}{\comment}{\endcomment}
\newenvironment{history}{\comment}{\endcomment}
\usepackage{cleveref}
\crefname{originallemma}{original lemma}{original lemmas}
\Crefname{originallemma}{Original lemma}{Original lemmas}
\makeatletter
\newlabel{lemma-c-square}{{0F16}{source}{}{sourcepointer-1}{}}
\newlabel{lemma-c-square@cref}{{[originallemma][0][]0F16}{[1][0][]source}{}{}{}}
\makeatother
\crefname{originallemma}{original lemma}{original lemmas}
\Crefname{originallemma}{Original lemma}{Original lemmas}
\makeatletter
\newlabel{lemma-two-out-of-three-blow-up-square}{{0EWP}{source}{}{sourcepointer-2}{}}
\newlabel{lemma-two-out-of-three-blow-up-square@cref}{{[originallemma][0][]0EWP}{[1][0][]source}{}{}{}}
\makeatother
\crefname{originalsection}{original section}{original sections}
\Crefname{originalsection}{Original section}{Original sections}
\makeatletter
\newlabel{section-compare}{{0EWK}{source}{}{sourcepointer-3}{}}
\newlabel{section-compare@cref}{{[originalsection][0][]0EWK}{[1][0][]source}{}{}{}}
\makeatother
\title{Descent of the blow-up square condition for complexes in the essential image of the derived pushforward}
\author{The Stacks Project Authors}
\date{stacks-project:sites-cohomology.tex:local-snapshot-20261011; GNU FDL 1.2 or later}
\begin{document}
\maketitle
\noindent\textit{Editorial note.} All delivered mathematical text is the author's own, copied from the pinned Stacks Project chapter file sites-cohomology.tex; no statement or proof step was rewritten, reordered or removed.\par\medskip
\noindent\textit{Editorial note.} The proof is the author's complete argument for both directions, including the induction-free vanishing step; every invoked input is named in the text.\par\medskip
\noindent\textit{Editorial note.} Every labelled reference occurring inside the counted spans is carried in assembly.external\_references as a non-load-bearing navigation pointer (same-chapter ones at their definition site, cross-chapter ones at the invocation inside the counted span), so no load-bearing mathematics is replaced by a stub. No bibliography citation occurs inside the counted spans.\par\medskip
\noindent\textit{Editorial note.} The checked chapter list for same-name collisions is the set of stacks chapters staged in this repository: chow, cohomology, divisors, etale-cohomology, flat, more-algebra, more-morphisms, morphisms, spaces-more-morphisms, varieties together with this card's chapter sites-cohomology.\par\medskip
\noindent\textit{Editorial note.} The preamble is the chapter's own retained minimal preamble (sources/stacks-sites-cohomology/retained-preamble.tex), a verbatim excerpt of the native preamble: the Xy-pic setup and the native theorem declarations with the mathematical macros.\par\medskip
\noindent\textit{Allowed standard prerequisites (editorial).}\begin{itemize}
\item The morphism of sites epsilon and its basic properties. Same chapter, Section section-compare, in particular epsilon\textasciicircum{}\{-1\} epsilon\_* = id. Supplies the adjunction and the identity used throughout the proof.
\item The canonical map attached to a blow-up square. Same chapter, Lemma lemma-c-square. Supplies the maps c\textasciicircum{}\{K'\}\_\{X,Z,Y,E\} whose being isomorphisms is characterized.
\item Two out of three for the blow-up square. Same chapter, Lemma lemma-two-out-of-three-blow-up-square. Shows that the third term of a distinguished triangle satisfies the condition when two terms do.
\end{itemize}\medskip
\ifdefined\hypertarget\hypertarget{sourcepointer-1}{}\fi
\noindent\textit{Source reference.} Original lemma 0F16: sources/stacks-sites-cohomology/pinned-original-source.tex, source lines 6001--6001. Navigation pointer invoked at proof line 6688; the counted span contains the pointer only and the cited result is not reproduced inside it, so the pointer is registered as a non-load-bearing navigation stub in the established form for Stacks material. No bibliography citation is involved.\par\medskip
\ifdefined\hypertarget\hypertarget{sourcepointer-2}{}\fi
\noindent\textit{Source reference.} Original lemma 0EWP: sources/stacks-sites-cohomology/pinned-original-source.tex, source lines 6046--6046. Navigation pointer invoked at proof line 6713; the counted span contains the pointer only and the cited result is not reproduced inside it, so the pointer is registered as a non-load-bearing navigation stub in the established form for Stacks material. No bibliography citation is involved.\par\medskip
\ifdefined\hypertarget\hypertarget{sourcepointer-3}{}\fi
\noindent\textit{Source reference.} Original section 0EWK: sources/stacks-sites-cohomology/pinned-original-source.tex, source lines 6177--6177. Navigation pointer invoked at proof line 6548; the counted span contains the pointer only and the cited result is not reproduced inside it, so the pointer is registered as a non-load-bearing navigation stub in the established form for Stacks material. No bibliography citation is involved.\par\medskip
\noindent
Let $\mathcal{C}$ be a category. Let
$\text{Cov}(\mathcal{C}) \supset \text{Cov}'(\mathcal{C})$
be two ways to endow $\mathcal{C}$ with the structure of a site.
Denote $\tau$ the topology corresponding to $\text{Cov}(\mathcal{C})$
and $\tau'$ the topology corresponding to $\text{Cov}'(\mathcal{C})$.
Then the identity functor on $\mathcal{C}$ defines a morphism
of sites
$$
\epsilon : \mathcal{C}_\tau \longrightarrow \mathcal{C}_{\tau'}
$$
where $\epsilon_*$ is the identity functor on underlying presheaves and
where $\epsilon^{-1}$ is the $\tau$-sheafification of a $\tau'$-sheaf.

\begin{lemma}
\label{lemma-descent-squares}
With $\epsilon : (\mathcal{C}_\tau, \mathcal{O}_\tau) \to
(\mathcal{C}_{\tau'}, \mathcal{O}_{\tau'})$ as above.
Let $A$ be a set and for $\alpha \in A$ let
$$
\xymatrix{
E_\alpha \ar[d] \ar[r] & Y_\alpha \ar[d] \\
Z_\alpha \ar[r] & X_\alpha
}
$$
be a commutative diagram in the category $\mathcal{C}$. Assume that
\begin{enumerate}
\item a $\tau'$-sheaf $\mathcal{F}'$ is a $\tau$-sheaf if
$\mathcal{F}'(X_\alpha) =
\mathcal{F}'(Z_\alpha) \times_{\mathcal{F}'(E_\alpha)} 
\mathcal{F}'(Y_\alpha)$ for all $\alpha$,
\item for $K'$ in $D(\mathcal{O}_{\tau'})$ in the essential image
of $R\epsilon_*$ the maps $c^{K'}_{X_\alpha, Z_\alpha, Y_\alpha, E_\alpha}$
of Lemma \ref{lemma-c-square}
are isomorphisms for all $\alpha$.
\end{enumerate}
Then $K' \in D^+(\mathcal{O}_{\tau'})$ is in
the essential image of $R\epsilon_*$ if and only if
the maps $c^{K'}_{X_\alpha, Z_\alpha, Y_\alpha, E_\alpha}$
are isomorphisms for all $\alpha$.
\end{lemma}

\begin{proof}
The ``only if'' direction is implied by assumption (2).
On the other hand, if $K'$ has a unique nonzero cohomology sheaf,
then the ``if'' direction follows from assumption (1).
In general we will use an induction argument to prove the
``if'' direction. Let us say an object $K'$ of $D^+(\mathcal{O}_{\tau'})$
satisfies (P) if the maps $c^{K'}_{X_\alpha, Z_\alpha, Y_\alpha, E_\alpha}$
are isomorphisms for all $\alpha \in A$.

\medskip\noindent
Namely, let $K'$ be an object of $D^+(\mathcal{O}_{\tau'})$
satisfying (P). Choose a distinguished triangle
$$
K' \to R\epsilon_*\epsilon^{-1}K' \to M' \to K'[1]
$$
in $D^+(\mathcal{O}_{\tau'})$ where the first arrow is the adjuntion map.
By (2) and Lemma \ref{lemma-two-out-of-three-blow-up-square}
we see that $M'$ has (P). On the other hand, applying
$\epsilon^{-1}$ and using that $\epsilon^{-1}R\epsilon_* = \text{id}$
by Section \ref{section-compare} we find that $\epsilon^{-1}M' = 0$.
In the next paragraph we will show $M' = 0$ which finishes the proof.

\medskip\noindent
Let $K'$ be an object of $D^+(\mathcal{O}_{\tau'})$
satisfying (P) with $\epsilon^{-1}K' = 0$. We will show $K' = 0$.
Namely, given $n \in \mathbf{Z}$ such that $H^i(K') = 0$
for $i < n$ we will show that $H^n(K') = 0$.
For $\alpha \in A$ we have a distinguished triangle
$$
R\Gamma_{\tau'}(X_\alpha, K') \to
R\Gamma_{\tau'}(Z_\alpha, K') \oplus
R\Gamma_{\tau'}(Y_\alpha, K') \to
R\Gamma_{\tau'}(E_\alpha, K') \to
R\Gamma_{\tau'}(X_\alpha, K')[1]
$$
by Lemma \ref{lemma-c-square}. Taking cohomology in degree $n$ and using
the assumed vanishing of cohomology sheaves of $K'$ we obtain an exact
sequence
$$
0 \to
H^n_{\tau'}(X_\alpha, K') \to
H^n_{\tau'}(Z_\alpha, K') \oplus
H^n_{\tau'}(Y_\alpha, K') \to
H^n_{\tau'}(E_\alpha, K')
$$
which is the same as the exact sequence
$$
0 \to
\Gamma(X_\alpha, H^n(K')) \to
\Gamma(Z_\alpha, H^n(K')) \oplus
\Gamma(Y_\alpha, H^n(K')) \to
\Gamma(E_\alpha, H^n(K'))
$$
We conclude that $H^n(K')$ is a a $\tau$-sheaf by assumption (1).
However, since the $\tau$-sheafification
$\epsilon^{-1}H^n(K') = H^n(\epsilon^{-1}K')$ is $0$ as
$\epsilon^{-1}K' = 0$
we conclude that $H^n(K') = 0$ as desired.
\end{proof}

\end{document}
